\section{Continuous Curvature}
\label{sec:curvature}

In this section, we introduce smooth and discrete notions of curvature. 
There are many ways to define curvature, each emphasizing different geometric insights and computational advantages. 
Some definitions are particularly intuitive and help develop geometric intuition, while others are better suited for efficient computation. 
Using these notions of curvature, we will state a version of the Gauss-Bonnet theorem for surfaces.



\subsection{Curvature of Surfaces}

Our goals for curvature are the following:
\begin{enumerate}
\item the plane has curvature zero
\item the unit sphere has curvature one
\item spheres with radius greater than one have curvature less than one and
spheres with radius less than one have curvature greater than one
\item curvature should be a signed quantity
\end{enumerate}

If the sphere has positive curvature, what does it mean to have negative curvature?
Consider standing on a saddle point, if we look in one direction the saddle is curving
one direction, say upwards, if we rotate $90^\circ$, the surface is curing the opposite direction, downward.  This is negative curvature.


For one dimensional curves in the plane, the osculating-circle is
met our criteria.
 Let us consider the \EMPH{osculating sphere} of a surface.
 Our criteria are met,
 but notice that at a saddle point on a surface it is not clear which sphere
best approximates the surface. See \figref{osculating-sphere} for two
equally reasonable ways to approximate a saddle with a sphere.
This suggests the osculating sphere is not ideal for defining the curvature
of a surface.

\begin{figure}[htb]
    \captionsetup[subfigure]{justification=centering}
    \centering
    \begin{subfigure}[b]{0.25\textwidth}
        \includegraphics[width=\textwidth]{curvature/sphere-above-saddle}
       \subcaption{}\label{fig:sphere-above-saddle}
    \end{subfigure}
        \hspace{1cm}
        \begin{subfigure}[b]{0.25\textwidth}
        \includegraphics[width=\textwidth]{curvature/sphere-below-saddle}
        \subcaption{}\label{fig:sphere-below-saddle}
        \end{subfigure}
    \caption{(\subref{fig:sphere-above-saddle}) The osculating sphere above the saddle.
        (\subref{fig:sphere-below-saddle}) The osculating sphere above the saddle.
    }
    \label{fig:osculating-sphere}
\end{figure}

Given a point on a smooth surface, we first compute the tangent plane at that point and then determine the corresponding normal vector. Any plane containing the normal vector is called a \emph{normal plane}. The intersection of the surface with a normal plane is a planar curve known as a \emph{normal section}; see \figref{normal-sections} for an illustration.


Each normal section has a curvature at the point of intersection. Among all such curvatures, let $(\kappa_1)$ denote the maximum curvature and $(\kappa_2)$ denote the minimum curvature. These quantities are called the \emph{principal curvatures} of the surface.
\begin{definition}[Gaussian Curvature One]
The \emph{Gaussian curvature} of a surface at \(p\) is

$$K(p)=\kappa_1\kappa_2.$$

\end{definition}
This definition agrees with the intuitive properties of curvature discussed above. In particular, Gaussian curvature is zero on locally flat surfaces, positive at points that bend in the same direction along both principal directions (such as a sphere), and negative at saddle points, where the surface bends in opposite directions along the principal directions.
Computing the normal vector is relatively simple in most cases, but
determining the orientation to maximize and minimize is often more difficult.



\begin{figure}[htb]
    \captionsetup[subfigure]{justification=centering}
    \centering
    \begin{subfigure}[b]{0.25\textwidth}
        \includegraphics[width=\textwidth]{curvature/normal-section-max}
       \subcaption{}\label{fig:normal-section-max}
    \end{subfigure}
        \hspace{1cm}
        \begin{subfigure}[b]{0.25\textwidth}
        \includegraphics[width=\textwidth]{curvature/normal-section-min}
        \subcaption{}\label{fig:normal-section-min}
        \end{subfigure}
    \caption{(\subref{fig:normal-section-max}) The normal section with maximum curvature.
        (\subref{fig:normal-section-min}) The normal section with minimum curvature.
    }
    \label{fig:normal-sections}
\end{figure}

A more computationally friendly definition is the following.
The two dimensional Gauss map provides a precise way to measure how a
surface curves in space.
Again, our surface $S$ is smooth, we have a well defined
tangent plane at any point $p\in S$  and a unit normal vector $N(p)$.
The \emph{Gauss map} is defined to be a map $g:S\rightarrow \Sp^2$
where, for $p\in S$ the value $g(p)$ is $N(p)$.


\begin{figure}[htb]
    \captionsetup[subfigure]{justification=centering}
    \centering
    \begin{subfigure}[b]{0.3\textwidth}
        \includegraphics[width=\textwidth]{curvature/gauss-map-convex}
       \subcaption{}\label{fig:gauss-map-convex}
    \end{subfigure}
        \hspace{1cm}
        \begin{subfigure}[b]{0.3\textwidth}
        \includegraphics[width=\textwidth]{curvature/gauss-map-nonconvex}
        \subcaption{}\label{fig:gauss-map-nonconvex}
        \end{subfigure}
    \caption{(\subref{fig:gauss-map-convex}) The Gauss map on a convex region.
        (\subref{fig:gauss-map-nonconvex}) The Gauss map on a nonconvex region.
    }
    \label{fig:normal-sections}
\end{figure}

The differential of the Gauss map, called the \emph{Weingarten map}, is a linear 
transformation that maps the tangent
space of the surface at $p$ to the tangent space of the sphere at $N(p)$.
$$dN_p:T_pS\rightarrow T_{N(p)}\Sp^2.$$
Since $N(p)$ is orthogonal to $T_pS,$ we naturally have
$$T_{N(p)}\cong T_pS.$$
The determinate of the Weingarten map represents how the normal
changes the area locally around $p$. If the normal changes rapidly, the surface 
curves sharply, if the normal changes slowly, the surface is nearly flat.

Let $U_1$ be a `small' set on a convex surface $S$. Where small means
contained in a disk whose intrinsic radius is approaching zero.
If the Gauss map is one-to-one on $U,$ and $g$ is orientation-preserving,
the curvature $K(U_1)$ of $U_1$ is defined to be the area of the image of $g(U_1)$
on $\Sp^2.$
Let $U_2$ be a nonconvex region on a surface on which $g$ is one-to-one and is orientation-reversing,
then $K(U_2)$ is defined as the negative of the area of $g(U_2)$ on $\Sp^2.$



We have a computable definition of the curvature of a surface. Let $W_p$ denote
the Weingarten map and let $k_1$ and $k_2$ be the eigenvalues.
\begin{definition}[Gaussian Curvature Two]
The \emph{Gaussian curvature} of a surface at \(p\) is

$$K(p)=\det(W_p).$$

\end{definition}
\noindent We see that this definition meets our basic criteria given at the beginning of the section.







\subsection{Computing Curvature}
We now discuss computing Gaussian curvature.
A \EMPH{quadratic form} is polynomial 
of the form $p(u,v)=c_1u^2+c_2uv+c_3v^2$ 
where $c_i\in R$.
Given a surface defined locally by $r(u,v)$,
we compute Gaussian curvature using a quadratic form called first fundamental form.

Let $E=r_u\cdot r_u, F=r_u\cdot r_v$ and  $G=r_v\cdot r_v$.
The \EMPH{first fundamental form}
is the quadratic form $\mathrm{I}=Edu^2+2Fdudv +Gdv^2$.
We summarize the first fundamental form as a matrix $$\mathrm{I}=\begin{bmatrix}
E & F \\
F & G 
\end{bmatrix}.$$
%We get a notion of length in the tangent space, an inner product on $Tp(S)$.
%If $x$ and $y$ are two tangent vectors
%then $$\mathrm{I}(x,y)=x^T\begin{bmatrix}
%E & F \\
%F & G 
%\end{bmatrix}y.$$

The first fundamental form enables us to  compute many interesting
things about our surface including arc length, angles between curves,
and the area of a subsurface.
For example,
given a `small' parallelogram $M$ on $S$ with corners $r(u,v),r(u+\epsilon u, v), r(u,v+\epsilon v)$ 
and $r(u+\epsilon u, v+\epsilon v)$ the rate of change of the area of $M$ in terms of $u$ and $v$ is 
$$dA=\sqrt{EG-F^2}dudv.$$
In \exref{stereo}, we will use the first fundamental form to compute the arc length of some curves.

A property of a surface is \emph{intrinsic} if it can be determined by working without referencing
an embedding of the surface in $\R^3.$ An intrinsic property can be computed by measuring angles,
distances and their derivatives. While we will not prove it, Gauss's Egregious (remarkable) Theorem states
\begin{theorem}[Theorem Egregious]
 Gaussian curvature and the first fundamental form intrinsic properties of a surface.
\end{theorem}

Often, we consider a surface, say the sphere, embedded in $\R^3.$ A property of a surface is
\emph{extrinsic} if it depends on how the surface sits in the ambient space. Extrinsic properties
tell us what the ambient space tells us about our surface. If you imagine you are a small
creature living on a surface, intrinsic properties are what you can compute. Extrinsic properties
are what an outside observer in $\R^3$ can compute.

The Brioschi formula, \cite{teo2024brioschi}, gives the Gaussian curvature using only the first fundamental form
\begin{equation}\label{eq:brioschi}
	K=\frac{\begin{vmatrix}
-\frac{1}{2}E_{vv}+F_{uv}-\frac{1}{2}G_{uu} & \frac{1}{2}E_u & F_u-\frac{1}{2}E_v\\
F_v-\frac{1}{2}G_u & E & F\\
\frac{1}{2}G_v & F & G
\end{vmatrix}-\begin{vmatrix}
0 & \frac{1}{2}E_v & \frac{1}{2}G_u\\
\frac{1}{2}E_v & E & F\\
\frac{1}{2}G_u & F & G
\end{vmatrix}}{(EG-F^2)^2}
\end{equation}
For most parameterized surfaces in the form $r(u,v),$ the partial derivatives used in the
first fundamental form are able to be computed relatively efficiently. Thus,
writing a computer program to compute the Gaussian curvature of a parameterized
surface at a point is straight forward.


While it is possible compute Gaussian curvature using only the first fundamental form,
 the second fundamental form makes the computation easier.
To this end, the unit normal vector at $p$ is given by $$n(p)=\frac{r_u\times r_v}{|r_u\times r_v|}.$$
Let $L=r_{uu}\cdot n, M=r_{uv}\cdot n$ and $N=r_{vv}\cdot n$ the
\EMPH{second fundamental form} is $\mathrm{I\!I}=Ldu^2+2Mdudv+Ndv^2$,
in matrix form $$\mathrm{I\!I}=\begin{bmatrix}
L & M \\
M & N 
\end{bmatrix}.$$
%Another inner product is given by $$\mathrm{I\!I}(x,y)=x^T\begin{bmatrix}
%L & M \\
%M & N 
%\end{bmatrix}y.$$
Combining the first and second fundamental forms we have
the \EMPH{Gaussian curvature} of a surface is
\begin{equation}\label{eqn:curve-dets}
 	K=\frac{\det(\mathrm{I\!I})}{\det(\mathrm{I})}.
\end{equation}


Let us look at an example to ensure that our definitions agree with our intuition.
\begin{example}[Curvature of the Sphere]\label{ex:compute-surface-curvature}
Consider the northern hemisphere of the unit sphere parameterized
by $$r(u,v)=(u,v,\sqrt{1-u^2-v^2})$$ and we wish to compute the Gaussian curvature at
the north pole $p=(0,0,1)$.
We have $$r_u|_p=(1,0,\frac{-u}{\sqrt{1-u^2-v^2}})|_p=(1,0,0)$$ and
$$r_v|_p=(0,1,\frac{-v}{\sqrt{1-u^2-v^2}})|_p=(0,1,0).$$
Then, $E=1, F=0,$ and $G=1$ so our first fundamental form is
$\mathrm{I}=\begin{bmatrix}
1 & 0 \\
0 & 1 
\end{bmatrix}.$
Our normal vector is $n(p)=(0,0,1)$ with 
$$r_{uu}|_p=(0,0,\frac{1-2u^2}{(1-u^2)^{\frac{3}{2}})})|_p=(0,0,1),
r_{uv}|_p=(0,0,0)$$ and 
$$r_{vv}|_p=(0,0,\frac{1-2v^2}{(1-v^2)^{\frac{3}{2}})})|_p=(0,0,1).$$
Thus, $L=1, M=0$ and $N=1$ and our second fundamental form is
$\mathrm{I\!I}=\begin{bmatrix}
1 & 0 \\
0 & 1 
\end{bmatrix}$
and our Gaussian curvature is $K=\frac{\det(\mathrm{I\!I})}{\det(\mathrm{I})}=\frac{1}{1}=1$ as expected.

\end{example}

To further illustrate the power of the fundamental forms, let us
compute
the arc length of circles on the sphere parallel 
to the $xy$ plane with fixed height $z=c$ for $-1<c<1$.
Stereographic projection, described in \secref{sphere-plane},
allows us map our problem to the plane and perform our computations there.

\begin{example}[Stereographic Projection Revisited\cite{christian-notes}]\label{ex:stereo}
Consider the two sphere with the north pole removed $\Sp^2 \setminus (0,0,1)$,
stereographic projection is a bijection between the points on $\Sp^2 \setminus (0,0,1)$ to the $\R^2$.
Consider a line from the north pole $(0,0,1)$ that intersects $(x,y,z)\in \Sp^2$ parametrized by 
$p(t)=(1-t)(0,0,1)+t(x,y,z)$. By considering the $z$ coordinate we determine the $t$ value where this line
intersects $\R^2$, namely $t=\frac{1}{1-z}.$
This gives the desired map shown in \figref{stereo} and in equation form
$$p(x,y,z)\to \left(\frac{x}{1-z},\frac{y}{1-z}\right).$$

\begin{figure}[htb]
	\centering
	\includegraphics[width=.3\textwidth]{curvature/stereo}
	\caption{A point on the sphere is mapped to a point on the plane by stereographic projection.}
	\label{fig:stereo}
\end{figure}
	
The inverse is given by $p^{-1}:\R^2\to \R^3$

	\begin{equation}\label{eqn:stereo}
		p^{-1}(u,v)=\left(\frac{2u}{u^2+v^2+1},\frac{2v}{u^2+v^2+1},\frac{u^2+v^2-1}{u^2+v^2+1}\right).	
	\end{equation}
To compute the first fundamental form of $p^{-1}(u,v)$ we take the partial derivatives

$$p^{-1}_u=\left(\frac{2v^2-2u^2+2}{(u^2+v^2+1)^2},\frac{-4uv}{(u^2+v^2+1)^2},\frac{4v}{(u^2+v^2+1)^2}\right)$$
and 
$$p^{-1}_v=\left(\frac{-4uv}{(u^2+v^2+1)^2},\frac{2v^2-2u^2+2}{(u^2+v^2+1)^2},\frac{4v}{(u^2+v^2+1)^2}\right).$$
Then, after some algebra,
$$E=p^{-1}_u\cdot p^{-1}_u=\frac{4}{(u^2+v^2+1)^2}$$
$$F=p^{-1}_v\cdot p^{-1}_v=\frac{4}{(u^2+v^2+1)^2}$$
and
$$M=p^{-1}_u\cdot p^{-1}_v=0.$$


Given a curve $\xi(t)=(u(t),v(t))$ on a regular surface, the arc
length of the curve
can be computed using the first fundamental form as follows.  Let
$s$ denote the arc length, then 
$$ds=\bigg | \frac{dr}{dt}\bigg | dt = \bigg | r_u\frac{du}{dt}+r_v\frac{dv}{dt}\bigg |dt
=\sqrt{(r_u^2 du^2+2r_ur_v du dv + r_v^2dv^2)}.$$

Using \eqnref{stereo}, our circle on the sphere maps 
to a circle in the plane because
	$$p^{-1}(u,v)=\frac{u^2+v^2-1}{u^2+v^2+1}=c$$
and we can compute the radius in terms of $c$
\begin{equation}\label{eqn:radius}
	u^2+v^2=\frac{1+c}{1-c}=k^2.
\end{equation}
Since we are in the plane, the circle $u^2+v^2=k^2$ can be parameterized
as $$\gamma(t)=(k\cos(t),k\sin(t))$$ with $0\leq t\leq 2\pi.$
So our curve becomes $p^{-1}\circ \gamma(t)$ on the sphere.
Computing the partial derivatives of $\gamma(t)$ gives
$$\gamma_u'=-k\sin(t)\hspace{1.3cm}  \gamma_v'=k\cos(t).$$
Now we use the first fundamental form

$$\int_{p^{-1}}\gamma ds=\int_{0}^{2\pi} ||(p^{-1}\circ \gamma)'(t)dt=\int_0^{2\pi}\sqrt{E(\gamma_u'(t))^2+2M\gamma_u'\gamma_v'+
F(\gamma_v'(t))^2}dt.$$
Substituting and simplifying using $E=F$ we obtain
$$\int_0^{2\pi}\frac{2k}{k^2+1}dt=\frac{4\pi k}{k^2+1}.$$
Simplifying further using \eqnref{radius}  our arc length is
$$2\pi\sqrt{1-c^2}.$$

Notice that this result agrees with the arc length determined by using
the Pythagorean theorem!

\end{example}


Although we will not prove it here, another interesting feature of the first fundamental form is 
that we can use it to determine if a map is conformal.
\begin{theorem}[First Fundamental Form and Conformal Maps]\label{thm:first-conformal}
	Let $S$ be a regular surface and $f:U\subset \R^2 \to V\subset S$ be a chart.
	Let $f$ have first fundamental form $\mathrm{I}=Edu^2+2Fdudv +Gdv^2$.
	Then $f$ is conformal if and only if $E=G$ and $F=0.$
\end{theorem}
%%%%%
\subsection{Gauss--Bonnet for Smooth Surfaces}

Now that we have defined the Euler characteristic and Gaussian curvature, we
have the tools to state a version of the Gauss--Bonnet theorem for smooth compact surfaces
without boundary. 

What makes this result particularly striking is that curvature is defined point-by-point using derivatives and local measurements, while the Euler characteristic depends only on the global structure of the surface. The Gauss--Bonnet theorem shows that when curvature is added up over an entire surface, the result depends only on the topology of the surface and not on its particular shape or embedding in space.

\begin{theorem}[Continuous Gauss-Bonnet Theorem - No Boundary] \label{thm:g-b-nb}

Let $S$ be a smooth, compact, orientable surface without boundary. Then
	$$\int_{S} K dA = 2\pi \chi(S)$$
	where  $K$ is Gaussian curvature and
	$\chi$ is the Euler characteristic.
\end{theorem}

Let us To illustrate the theorem, consider the unit sphere $\Sp^2$. The Gaussian curvature of the unit sphere is constant and equal to,K=1. Since the surface area of the unit sphere is $4\pi$,
$$\int_{S^2} dA=4\pi.$$
On the other hand, the Euler characteristic of the sphere is
$$\chi(S^2)=2$$
multiply by the stated $2\pi$ and the theorem holds for $\Sp^2.$

We require our surface to be orientable because we rely on a consistent
normal vector to compute curvature.  We will discuss a version for non-orientable surfaces.
There are versions of the Gauss--Bonnet theorem for non orientable surfaces.
We require out surface to be compact because we are defining the Euler characteristic
in terms of triangulation. For example $\R^2$ is not compact, we need an infinite number of
triangles to triangulate it.
There are versions of the Gauss--Bonnet theorem for noncompact surfaces.

In \secref{intro}, we stated the \thmref{simple-bonnet} the theorem of turning tangents for simple
polygons. Simple polygons in the plane have boundaries, with exterior angles that are contributing
the curvature.
Our interpretation of the curvature of a plane curve is how quickly the unit
tangent vector is changing as we move along the curve.
On a surface, the rate of change of the unit tangent vector depends on two factors
\begin{enumerate}
\item the curvature of the surface itself and
\item how the curve is bending within the surface.
\end{enumerate}
We decompose curvature in to tow components, one normal to the surface
and one tangent to the surface.
The tangential component measure how much the curve bends within the surface
and is called \emph{geodesic curvature}.
\begin{definition}[Geodesic Curvature]
Let $S$ be a smooth surface and let
$\gamma:[a,b]\to S$
be a unit-speed curve on $S$. Let $T(s)=\gamma'(s)$ denote the unit
tangent vector to the curve, and let $N(s)$ denote the unit normal vector
to the surface along $\gamma$.
Since $\kappa(s)\mathbf{n}(s)=T'(s),$
the curvature vector $T'(s)$ can be decomposed into components tangent and
normal to the surface:
$$T'(s)=\kappa_g(s)\,\mathbf{u}(s)+\kappa_n(s)\,N(s),$$
where $\mathbf{u}(s)$ is a unit vector tangent to the surface and
orthogonal to $T(s)$.
The scalar $\kappa_g(s)$ is called the \emph{geodesic curvature} of
$\gamma$ at $s$.
\end{definition}


We compute the geodesic curvature on a parameterized as follows.
First, find the tangent plane to the surface at the point and the unit normal vector.
Let $U:\mathcal{U}\to \RR^3$ be a parameterized chart on a surface $S$ with vector $n(u,v)$ normal
to the surface
and let $\gamma(s)$ be a curve in $U$.
Then $\mathbf{v}=n(\gamma(s))\times T$ is in the tangent plane of the surface since
it is perpendicular to $n$. Moreover, $\mathbf{v}(s)$ is normal to $\gamma$ 
from the perspective of someone living on the surface. 
The geodesic curvature tells us the rate of change of $\gamma'$ with respect 
to $\mathbf{v}(s)$.
Then, the geodesic curvature is given by 
\begin{equation} \label{eqn:geodesic}
	k_g=\gamma'' \cdot (\mathbf{v}\times \gamma').
\end{equation}

On a surface, a \EMPH{geodesic} is a curve that is a shortest path
between two points in the surface. 
For example, on $\Sp^2$ great circles are geodesic.
Intuitively, the geodesic curvature of a one dimensional curve on a surface
is the curvature as it would be seen from someone living on the surface.

\begin{example}[Circles on the Sphere]\label{eqn:circles-on-sphere}
	We can rotate the sphere so that the circle is parallel to the $xy$-plane.
	Then the radius of the circle $r$ is related to the height $h$ of the circle above the $xy$-plane
	by $h^2+r^2=1$. See \figref{geodesic} for an example.
	A parameterization by arc length is given by
	$$\gamma(\theta)=\left(r\cos\left(\frac{\theta}{r}\right),r\sin\left(\frac{\theta}{r}\right),\sqrt{1-r^2}\right),$$
	so
	$$\gamma'(\theta)=\left(-\sin\left(\frac{\theta}{r}\right),\cos\left(\frac{\theta}{r}\right),0\right)$$
	and
	$$\gamma''(\theta)=\left(-\frac{1}{r}\cos\left(\frac{\theta}{r}\right),-\frac{1}{r}\sin\left(\frac{\theta}{r}\right),0\right).$$
	Since we are on the sphere, the normal vector is equal to the point on the sphere
	so $$n(\theta)=\left(r\cos\left(\frac{\theta}{r}\right),r\sin\left(\frac{\theta}{r}\right),\sqrt{1-r^2}\right).$$
	Then 
	$$k_g=\gamma''\cdot V(\theta)=\frac{\sqrt{1-r^2}}{r}.$$
	Notice that if $h=0$ in \figref{geodesic}, then $r=1$ and the geodesic curvature is zero which
	agrees with our intuition since we have a great circle.
\end{example}

\begin{figure}[htb]
	\centering
	\includegraphics[width=.25\textwidth]{curvature/geodesic}
	\caption{Computing the geodesic curvature.}
	\label{fig:geodesic}
\end{figure}



For smooth surfaces with boundary, even if the boundary that has corners
the above theorem can be adapted to include the curvature on the boundary
to obtain the following.

\begin{theorem}[Continuous Gauss-Bonnet with Boundary] \label{thm:g-b-bd}
Let $S$ be a smooth, compact, orientable surface with boundary.
Then
	$$\int_{S} K dA+ \int_{\partial S} k_g ds + \sum_i \beta_i= 2\pi \chi(S)$$
	where  $K$ is Gaussian curvature,
	 $k_g$ is the geodesic curvature,
	 each $\beta_i$  is an exterior angle at a vertex of the boundary and
	$\chi$ is the Euler characteristic.
\end{theorem}

To illustrate the Gauss--Bonnet theorem with boundary, consider the unit disk
$$D={(x,y)\in \mathbb{R}^2 : x^2+y^2\leq 1}.$$
Since the disk lies in the plane, its Gaussian curvature is zero.

The boundary of the disk is the unit circle $\partial D=\Sp^1$. The geodesic curvature of the unit circle is constant and equal to one, while its length is $2\pi$. Hence

$$\int_{\partial D} k_gds = \int_{\partial D} 1ds=2\pi.$$

On the other hand, the disk is contractible and therefore has Euler characteristic
$\chi(D)=1.$
Multiplying by the factor $2\pi$ appearing in the theorem gives
$2\pi\chi(D)=2\pi$
which verifies the Gauss--Bonnet theorem for the disk.

This example is particularly illuminating because all of the curvature is concentrated on the boundary. Although the disk itself is perfectly flat, the curvature contributed by its boundary is exactly what is needed to recover the Euler characteristic of the surface.


\noindent \textbf{Exercises}


\begin{enumerate}
	\item Given the curve $\gamma(t)=(t,t^2)$ for $t\in[-1,1]$ compute 
	the radius of the osculating sphere at $t=0.$
	
	\item Use any resource to convince yourself that the Weingarten has real eigenvalues.
	
	\item 
	
\end{enumerate}

\pagebreak